\documentclass[11pt,a4paper]{article}
\usepackage[T2A]{fontenc}
\usepackage[utf8]{inputenc}
\usepackage[russian]{babel}
\usepackage{geometry}
\geometry{left=19mm,right=19mm,top=20mm,bottom=20mm}
\usepackage{amsmath,amssymb,mathtools}
\usepackage{array,booktabs,longtable}
\usepackage{xcolor,hyperref,tikz}
\usetikzlibrary{arrows.meta,calc,angles,quotes}
\hypersetup{colorlinks=true,linkcolor=blue!55!black,urlcolor=blue!55!black}
\setlength{\parindent}{0pt}
\setlength{\parskip}{5pt}
\setlength{\emergencystretch}{3em}
\sloppy
\newcommand{\task}[2]{\newpage\phantomsection\label{task:#1}\section*{Задание №#1. #2}}
\newcommand{\err}[1]{\par\textcolor{red}{\textbf{Ошибка:}} #1\par}
\newcommand{\keyidea}[1]{\par\textcolor{blue!60!black}{\textbf{Ключевая идея:}} #1\par}
\newcommand{\conclusion}[1]{\par\textcolor{teal!60!black}{\textbf{Итог:}} #1\par}
\newcommand{\itemhead}[1]{\par\medskip\textbf{#1}\par}
\newcommand{\vishead}[1]{\par\medskip{\color{blue!55!black}\textbf{Наглядная схема: #1}}\par\smallskip}
\newcommand{\viscaption}[1]{\par\smallskip{\small\color{black!65}#1}\par\medskip}
\begin{document}
\thispagestyle{empty}
\vspace*{49mm}
\begin{center}
{\Huge\bfseries Ревью домашней работы\par}
\vspace{17mm}
{\Large Автор: Лёвин Артём Александрович\par}
\vspace{7mm}
{\large 24 октября 2026 года\par}
\end{center}
\vfill
\begin{center}
\begin{tikzpicture}
\draw[blue!40!black,line width=.8pt] (-4,0) .. controls (-2.2,.7) and (-1,-.55) .. (0,0)
 .. controls (1,.55) and (2.2,-.7) .. (4,0);
\end{tikzpicture}
\end{center}
\newpage
\section*{Сводная таблица ошибок и незавершённых решений}
\small
\setlength{\tabcolsep}{3pt}
\begin{longtable}{@{}p{.95cm}p{3.1cm}p{5.3cm}p{2.8cm}p{3.2cm}@{}}
\toprule
\textbf{№} & \textbf{Тема} & \textbf{Суть ошибки} & \textbf{Ваш ответ} & \textbf{Правильный ответ}\\
\midrule
\endfirsthead
\toprule
\textbf{№} & \textbf{Тема} & \textbf{Суть ошибки} & \textbf{Ваш ответ} & \textbf{Правильный ответ}\\
\midrule
\endhead
\hyperref[task:1]{1} & Описанная трапеция & Решение и ответ не представлены. & Не указан & $\dfrac72$\\
\hyperref[task:3]{3} & Прямая призма & Не вычислена площадь поверхности. & Не указан & $216$\\
\hyperref[task:5]{5} & Вероятность попаданий & Не сравнены два варианта стрельбы. & Не указан & $0{,}48$\\
\hyperref[task:6]{6} & Математическое ожидание & Не составлено уравнение для цены полиса. & Не указан & $3000$ руб.\\
\hyperref[task:9]{9} & Производная по графику & Найден наклон не той прямой; потерян минус. & $1{,}5$ & $-0{,}5$\\
\hyperref[task:12]{12} & Квадратичная функция & Ошибки при нахождении $a$ и подстановке $-5$. & $-0{,}6$ & $12$\\
\hyperref[task:14]{14} & Отбор корней на отрезке & Пункт (а) верен; не выполнен пункт (б). & Общие серии верны; пункт (б) отсутствует & $2\pi,\ \dfrac{10\pi}{3},\ \dfrac{14\pi}{3}$\\
\hyperref[task:16]{16} & Показательное неравенство & Вместо $152$ переписано $159$; решение оборвано. & Не указан & $(-3;0)\cup(0;\frac32)$\\
\hyperref[task:17]{17} & Минимум функции & Перепутан порядок членов в формуле производной; минимум не обоснован. & $180$ м & $180$ м\\
\bottomrule
\end{longtable}
\normalsize
\newpage
\thispagestyle{empty}
\vspace*{\fill}
\begin{center}{\Huge\bfseries Разбор ошибок}\end{center}
\vspace*{\fill}

\task{1}{Окружность, вписанная в трапецию}
\itemhead{Текст задания}
Периметр прямоугольной трапеции, описанной около окружности, равен $32$.
Найдите радиус окружности, если большая боковая сторона равна $9$.
\itemhead{Ваше решение}
Письменное решение на представленной странице отсутствует.
\itemhead{Ваш ответ}
Не указан.
\err{Задание не завершено: не использованы свойства трапеции, описанной около окружности.}
\itemhead{Где именно возникает проблема}
После условия нет первого математического шага. Последнего правильного преобразования и первого неправильного преобразования здесь нет: работа не начата.
\itemhead{Почему без этого нельзя получить ответ}
Периметр сам по себе не задаёт высоту произвольной трапеции. Нужно воспользоваться дополнительным условием касания всех четырёх сторон окружности.
У описанного четырёхугольника суммы противоположных сторон равны. Кроме того, в прямоугольной трапеции одна боковая сторона перпендикулярна основаниям и является высотой; наклонная боковая сторона не короче высоты.
\keyidea{Сумма оснований равна сумме боковых сторон. Из периметра можно определить их сумму, затем найти высоту и радиус.}

\vishead{прямоугольная трапеция и вписанная окружность}
\begin{center}
\begin{tikzpicture}[scale=.84,line cap=round]
  % Размеры строго согласованы: высота 7, наклонная боковая сторона 9,
  % сумма оснований 16 (единицы рисунка масштабированы в 5/7 раза).
  \coordinate (A) at (0,0);
  \coordinate (B) at (7.7346,0);
  \coordinate (C) at (3.6940,5);
  \coordinate (D) at (0,5);
  \fill[blue!3] (A)--(B)--(C)--(D)--cycle;
  \draw[thick,blue!65!black] (A)--(B)--(C)--(D)--cycle;
  \draw[teal!70!black,thick] (2.5,2.5) circle (2.5);
  \fill[teal!70!black] (2.5,2.5) circle (1.4pt);
  \draw[teal!65!black,dashed] (2.5,2.5)--(2.5,0);
  \draw[teal!65!black,-{Stealth[length=2.5mm]}] (2.5,2.5)--(4.27,4.27)
    node[pos=.55,below left] {$r=3{,}5$};
  \draw (0,.35)--(.35,.35)--(.35,0);
  \node[below left] at (A) {$A$};
  \node[below right] at (B) {$B$};
  \node[above right] at (C) {$C$};
  \node[above left] at (D) {$D$};
  \node[left] at (0,2.5) {$h=7$};
  \node[right] at (5.7,2.65) {$9$};
  \node[below] at (3.9,0) {$a$};
  \node[above] at (1.8,5) {$b$};
\end{tikzpicture}
\viscaption{Поскольку окружность касается обоих оснований, её диаметр равен высоте трапеции: $2r=h$.}
\end{center}

\itemhead{Исправление}
Обозначим основания через $a,b$, боковые стороны через $c,d$.
Начать следует с $a+b=c+d$ и $a+b+c+d=32$.
\itemhead{Полное правильное решение}
Так как трапеция описана около окружности,
\[
a+b=c+d.
\]
По условию периметр равен $32$, поэтому
\[
(a+b)+(c+d)=32,\qquad 2(c+d)=32,\qquad c+d=16.
\]
Большая боковая сторона равна $9$; другая равна
\[
16-9=7.
\]
В прямоугольной трапеции перпендикулярная боковая сторона равна высоте. Наклонная боковая сторона --- гипотенуза прямоугольного треугольника, а потому она не меньше высоты. Значит, высота $h=7$.
Вписанная окружность касается обоих параллельных оснований. Расстояние между этими прямыми равно диаметру окружности:
\[
h=2r,\qquad r=\frac{7}{2}=3{,}5.
\]
\itemhead{Проверка}
Сумма боковых сторон действительно равна $9+7=16$, а сумма оснований равна $16$ по свойству описанного четырёхугольника. Тогда периметр равен $16+16=32$. Диаметр окружности $2\cdot 3{,}5=7$ совпадает с высотой.
\conclusion{Радиус окружности равен $3{,}5$.}
\itemhead{Задача на закрепление}
Периметр прямоугольной трапеции, описанной около окружности, равен $40$. Большая боковая сторона равна $12$. Найдите радиус окружности.

\task{3}{Площадь поверхности прямой призмы}
\itemhead{Текст задания}
Основанием прямой треугольной призмы является прямоугольный треугольник с катетами $3\sqrt3$ и $4\sqrt3$ и гипотенузой, равной боковому ребру призмы. Найдите площадь поверхности призмы.
\itemhead{Ваше решение}
Решение на представленной странице не записано.
\itemhead{Ваш ответ}
Не указан.
\err{Не найдены гипотенуза основания, высота призмы и площадь её полной поверхности.}
\itemhead{Где именно возникает проблема}
Решение прерывается до вычислений. Нет корректного начального шага, от которого можно продолжить записанный ход.
\itemhead{Почему это важно}
Полная поверхность прямой призмы состоит из двух равных оснований и боковой поверхности. Если посчитать только площадь треугольника, боковые грани останутся неучтёнными. Чтобы найти боковую поверхность, нужны периметр основания и высота призмы.
\keyidea{Использовать теорему Пифагора и формулу $S_{\text{полн}}=2S_{\text{осн}}+P_{\text{осн}}h$.}

\vishead{прямая треугольная призма}
\begin{center}
\begin{tikzpicture}[scale=.95,line join=round]
  \coordinate (A) at (0,0); \coordinate (B) at (4.4,0);
  \coordinate (C) at (1.65,1.55);
  \coordinate (A1) at (0,3.15); \coordinate (B1) at (4.4,3.15);
  \coordinate (C1) at (1.65,4.70);
  \fill[blue!6] (A)--(B)--(B1)--(A1)--cycle;
  \fill[teal!8] (A1)--(B1)--(C1)--cycle;
  \draw[blue!70!black,thick] (A)--(B)--(B1)--(A1)--cycle;
  \draw[blue!70!black,thick] (A1)--(C1)--(B1);
  \draw[blue!70!black,thick] (B)--(C)--(C1);
  \draw[blue!70!black,dashed] (A)--(C)--(C1);
  \draw[blue!70!black,dashed] (A)--(A1);
  \node[below left] at (A) {$A$}; \node[below right] at (B) {$B$};
  \node[right] at (C) {$C$};
  \node[left] at (A1) {$A_1$}; \node[right] at (B1) {$B_1$};
  \node[above] at (C1) {$C_1$};
  \node[below] at (2.1,0) {$3\sqrt3$};
  \node[above left] at (.75,.82) {$4\sqrt3$};
  \node[right] at (3.3,.88) {$5\sqrt3$};
  \node[left] at (-.1,1.58) {$h=5\sqrt3$};
  \node[font=\scriptsize,teal!60!black] at (.8,.2) {$\angle BAC=90^\circ$};
\end{tikzpicture}
\viscaption{Трёхмерная схема дана в проекции: треугольник основания прямоугольный, боковые рёбра перпендикулярны основаниям.}
\end{center}

\itemhead{Исправление}
Сначала определить гипотенузу прямоугольного треугольника. По условию именно она равна высоте призмы.
\itemhead{Полное правильное решение}
Пусть катеты основания равны $a=3\sqrt3$ и $b=4\sqrt3$.
По теореме Пифагора гипотенуза $c$ удовлетворяет условию
\[
c^2=(3\sqrt3)^2+(4\sqrt3)^2=27+48=75,
\quad c=\sqrt{75}=5\sqrt3.
\]
Так как призма прямая и её боковое ребро равно гипотенузе основания,
\[
h=5\sqrt3.
\]
Площадь одного основания:
\[
S_{\text{осн}}=\frac12 ab=\frac12\cdot3\sqrt3\cdot4\sqrt3=18.
\]
Периметр основания:
\[
P_{\text{осн}}=3\sqrt3+4\sqrt3+5\sqrt3=12\sqrt3.
\]
Боковая поверхность прямой призмы:
\[
S_{\text{бок}}=P_{\text{осн}}h=12\sqrt3\cdot5\sqrt3=180.
\]
Полная поверхность:
\[
S_{\text{полн}}=2\cdot18+180=216.
\]
\itemhead{Проверка}
Два треугольных основания дают площадь $36$. Три прямоугольные боковые грани дают площади $45$, $60$ и $75$; их сумма $180$. Получаем $36+180=216$.
\conclusion{Площадь поверхности призмы равна $216$.}
\itemhead{Задача на закрепление}
Основание прямой треугольной призмы --- прямоугольный треугольник с катетами $6$ и $8$. Боковое ребро равно гипотенузе основания. Найдите площадь полной поверхности призмы.

\task{5}{Сравнение вероятностей попадания}
\itemhead{Текст задания}
Стрелок попадает в первую мишень первым выстрелом с вероятностью $0{,}5$, вторым --- с вероятностью $0{,}6$. Во вторую мишень первым выстрелом он попадает с вероятностью $0{,}6$, вторым --- с вероятностью $0{,}7$. Можно сделать два выстрела по первой и один по второй мишени либо один по первой и два по второй. Выберите вариант с большей вероятностью попасть в обе мишени и укажите эту вероятность.
\itemhead{Ваше решение}
Расчёты на представленной странице отсутствуют.
\itemhead{Ваш ответ}
Не указан.
\err{Не рассмотрены два допустимых варианта и не найдена вероятность попадания в обе мишени.}
\itemhead{Где именно возникает проблема}
Нет первого шага сравнения вариантов; ошибочных промежуточных вычислений в записи нет.
\itemhead{Почему простого сложения недостаточно}
Для мишени, в которую стреляют дважды, событие «попасть хотя бы один раз» включает три возможных сочетания результатов. Удобнее рассмотреть противоположное событие: оба раза промахнуться. Затем вероятность поражения обеих мишеней при независимых выстрелах получают умножением вероятностей соответствующих попаданий.
\keyidea{Для двух выстрелов вероятность хотя бы одного попадания равна единице минус вероятность двух промахов. Затем сравнить произведения для обоих вариантов.}

\vishead{сравнение двух вариантов стрельбы}
{\small\setlength{\tabcolsep}{4pt}
\begin{center}
\begin{tabular}{@{}p{38mm}p{42mm}p{39mm}p{29mm}@{}}
\toprule
\textbf{Вариант} & \textbf{Мишень с двумя выстрелами} & \textbf{Мишень с одним выстрелом} & \textbf{Обе мишени}\\
\midrule
Два по первой, один по второй & $1-0{,}5\cdot0{,}4=0{,}8$ & $0{,}6$ & $\boldsymbol{0{,}48}$\\
Один по первой, два по второй & $1-0{,}4\cdot0{,}3=0{,}88$ & $0{,}5$ & $0{,}44$\\
\bottomrule
\end{tabular}
\end{center}}
\viscaption{Для двух выстрелов сначала считаем вероятность хотя бы одного попадания. Итог по двум мишеням --- произведение.}

\itemhead{Исправление}
Записать отдельно вероятность успеха для варианта с двумя выстрелами по первой мишени и варианта с двумя выстрелами по второй мишени.
\itemhead{Полное правильное решение}
\textbf{Первый вариант:} два выстрела по первой мишени и один по второй.
Вероятности промахов по первой мишени равны $1-0{,}5=0{,}5$ и $1-0{,}6=0{,}4$.
При независимых выстрелах вероятность двух промахов равна
\[
0{,}5\cdot0{,}4=0{,}20.
\]
Поэтому вероятность хотя бы одного попадания в первую мишень:
\[
1-0{,}20=0{,}80.
\]
Вероятность попадания во вторую мишень одним выстрелом равна $0{,}6$.
Следовательно,
\[
P_1=0{,}80\cdot0{,}6=0{,}48.
\]
\textbf{Второй вариант:} один выстрел по первой мишени и два по второй.
Вероятности двух промахов по второй мишени равны $0{,}4$ и $0{,}3$; их произведение:
\[
0{,}4\cdot0{,}3=0{,}12.
\]
Значит, вероятность попасть во вторую мишень хотя бы раз равна
\[
1-0{,}12=0{,}88.
\]
Для первой мишени вероятность попадания одним выстрелом $0{,}5$. Получаем
\[
P_2=0{,}5\cdot0{,}88=0{,}44.
\]
Сравнение показывает: $0{,}48>0{,}44$.
\itemhead{Проверка}
Обе вероятности находятся между нулём и единицей. Разность $0{,}48-0{,}44=0{,}04>0$, следовательно, первый вариант лучше.
\conclusion{Следует выбрать два выстрела по первой мишени и один по второй; вероятность успеха $0{,}48$.}
\itemhead{Задача на закрепление}
В первую мишень стрелок попадает первым выстрелом с вероятностью $0{,}4$, вторым --- с вероятностью $0{,}8$; во вторую --- с вероятностями $0{,}7$ и $0{,}6$. Сравните варианты «два выстрела по первой и один по второй» и «один по первой и два по второй» и найдите наибольшую вероятность поразить обе мишени.

\task{6}{Математическое ожидание дохода}
\itemhead{Текст задания}
При наступлении страхового случая страховая компания выплачивает $100$ тысяч рублей. Вероятность выплаты по одному полису равна $0{,}01$. Математическое ожидание дохода компании от продажи одного полиса равно $2000$ рублей. Найдите стоимость полиса в рублях.
\itemhead{Ваше решение}
Решение на представленной странице не записано.
\itemhead{Ваш ответ}
Не указан.
\err{Не выражен доход компании с учётом возможной страховой выплаты.}
\itemhead{Где именно возникает проблема}
Отсутствует математическая модель. Ни верных, ни ошибочных промежуточных шагов не представлено.
\itemhead{Почему нужно учесть два исхода}
Страховая компания получает стоимость полиса в любом случае, но только при страховом случае дополнительно выплачивает клиенту $100000$ рублей. Значит, доход при двух исходах отличается; математическое ожидание является взвешенным средним этих доходов.
\keyidea{Ожидаемый доход равен цене полиса за вычетом ожидаемой страховой выплаты.}

\vishead{исходы страхового договора}
{\small\setlength{\tabcolsep}{4pt}
\begin{center}
\begin{tabular}{@{}p{45mm}p{24mm}p{38mm}p{43mm}@{}}
\toprule
\textbf{Исход} & \textbf{Вероятность} & \textbf{Доход компании} & \textbf{Вклад в ожидание}\\
\midrule
Выплаты нет & $0{,}99$ & $C$ & $0{,}99C$\\
Выплата есть & $0{,}01$ & $C-100000$ & $0{,}01(C-100000)$\\
\midrule
\multicolumn{3}{@{}l}{\textbf{Сумма вкладов}} & $\boldsymbol{C-1000}$\\
\bottomrule
\end{tabular}
\end{center}}
\viscaption{Стоимость полиса $C$ поступает при любом исходе; страховая выплата возможна только во втором случае.}

\itemhead{Исправление}
Обозначить стоимость полиса через $C$ и составить уравнение ожидания дохода.
\itemhead{Полное правильное решение}
Пусть полис стоит $C$ рублей.
Вероятность страхового случая равна $0{,}01$, а его отсутствия --- $0{,}99$.
Если выплаты нет, доход равен $C$. Если выплата есть, доход равен $C-100000$.
Тогда математическое ожидание:
\[
M=0{,}99C+0{,}01(C-100000).
\]
Раскрываем скобки и приводим подобные слагаемые:
\[
M=0{,}99C+0{,}01C-1000=C-1000.
\]
По условию $M=2000$, поэтому
\[
C-1000=2000,\qquad C=3000.
\]
\itemhead{Проверка}
Ожидаемая выплата равна $0{,}01\cdot100000=1000$ рублей. При цене полиса $3000$ рублей ожидаемый доход компании составляет $3000-1000=2000$ рублей, как и требуется.
\conclusion{Стоимость полиса равна $3000$ рублей.}
\itemhead{Задача на закрепление}
Страховая компания выплачивает $80000$ рублей при страховом случае с вероятностью $0{,}025$. Математическое ожидание дохода от одного проданного полиса должно равняться $1500$ рублей. Найдите стоимость полиса.

\task{9}{Производная как наклон касательной}
\itemhead{Текст задания}
На рисунке изображены график функции $y=f(x)$ и касательная к нему в точке с абсциссой $x_0$. Найдите значение $f'(x_0)$.
\itemhead{Ваше решение}
Записано: $\tg\alpha=\frac32=1{,}5$. На рисунке построен вспомогательный треугольник с возрастающей наклонной стороной, отличной от касательной.
\itemhead{Ваш ответ}
$1{,}5$.
\err{Найдён наклон возрастающей вспомогательной прямой вместо убывающей касательной.}
\itemhead{Где именно возникает ошибка}
Последнее верное наблюдение: производная в точке определяется угловым коэффициентом касательной. Первый неверный шаг --- использование отношения $3/2$ для наклона линии, которая не совпадает с касательной на рисунке. Поэтому получен положительный знак, хотя касательная убывает.
\itemhead{Почему этот шаг неверен}
Производная $f'(x_0)$ равна наклону \emph{касательной}, а не наклону произвольно проведённого отрезка и не наклону хорды графика. Когда при движении вправо вдоль касательной значение $y$ уменьшается, приращение ординаты отрицательно. Следовательно, отношение приращения ординаты к приращению абсциссы тоже отрицательно.
\keyidea{Выбирать две узловые точки именно на нарисованной касательной и вычислять отношение вертикального изменения к горизонтальному.}
\itemhead{Исправление}
По клеткам касательная проходит, например, через точки $(-3;-2)$ и $(-1;-3)$. При движении вправо на $2$ клетки она опускается на $1$ клетку.

\vishead{отрицательный наклон касательной}
\begin{center}
\begin{tikzpicture}[x=.85cm,y=.72cm]
 \draw[step=1,gray!23,very thin] (-5,-5) grid (3,1);
 \draw[-{Stealth},gray!75] (-5.35,0)--(3.35,0) node[right] {$x$};
 \draw[-{Stealth},gray!75] (0,-5.25)--(0,1.4) node[above] {$y$};
 \draw[blue!70!black,line width=1.15pt] (-5,-1)--(3,-5);
 \fill[blue!70!black] (-3,-2) circle (2.2pt) node[above left] {$A(-3;-2)$};
 \fill[blue!70!black] (-1,-3) circle (2.2pt) node[below right] {$B(-1;-3)$};
 \draw[teal!75!black,dashed,thick] (-3,-2)--(-1,-2)--(-1,-3);
 \node[above,teal!75!black] at (-2,-2) {$\Delta x=2$};
 \node[right,teal!75!black] at (-1,-2.52) {$\Delta y=-1$};
 \node[font=\small,blue!70!black] at (1.4,-4.35) {касательная};
\end{tikzpicture}
\viscaption{Опорные точки лежат на касательной. При перемещении вправо на $2$ клетки ордината уменьшается на $1$ клетку: $k=\Delta y/\Delta x=-1/2$.}
\end{center}

\itemhead{Полное правильное решение}
Возьмём две точки касательной:
\[
A(-3;-2),\qquad B(-1;-3).
\]
Её угловой коэффициент равен
\[
k=\frac{y_B-y_A}{x_B-x_A}=
\frac{-3-(-2)}{-1-(-3)}=\frac{-1}{2}=-0{,}5.
\]
Касательная проведена в точке с абсциссой $x_0$, поэтому
\[
f'(x_0)=k=-0{,}5.
\]
\itemhead{Проверка}
Нарисованная касательная идёт сверху слева вниз направо, следовательно, её угловой коэффициент отрицателен. Значение $-0{,}5$ согласуется с рисунком, а $1{,}5$ --- нет.
\conclusion{Значение производной равно $-0{,}5$.}
\itemhead{Задача на закрепление}
Касательная к графику функции в точке с абсциссой $x_0$ проходит через точки $(-2;4)$ и $(4;1)$. Найдите $f'(x_0)$.

\task{12}{Квадратичная функция по двум точкам}
\itemhead{Текст задания}
На рисунке изображён график функции $f(x)=ax^2-0{,}8x+c$. Найдите $f(-5)$.
\itemhead{Ваше решение}
По точкам $(0;-2)$ и $(2;-2)$ правильно определено $c=-2$, но затем записано $a=0{,}04$. После этого выполнена подстановка:
\[
f(-5)=0{,}04\cdot25+0{,}40-2=-0{,}6.
\]
\itemhead{Ваш ответ}
$-0{,}6$.
\err{При решении уравнения для коэффициента $a$ потерян десятичный разряд; при подстановке $-5$ допущена ещё одна ошибка.}
\itemhead{Где именно возникают ошибки}
Последний правильный шаг: из $f(0)=-2$ получено $c=-2$.
Первый ошибочный шаг: из уравнения $4a=1{,}6$ записано $a=0{,}04$ вместо $a=0{,}4$.
Второй неверный переход: произведение $(-0{,}8)\cdot(-5)$ заменено на $0{,}40$, хотя оно равно $4$.
\itemhead{Почему это неверно}
Деление $1{,}6$ на $4$ даёт $0{,}4$, поскольку $4\cdot0{,}4=1{,}6$. Результат $0{,}04$ в десять раз меньше правильного.
Кроме того, произведение двух отрицательных чисел положительно, а $0{,}8\cdot5=4$. Нельзя изменять разрядность результата при умножении.
\keyidea{Считать коэффициенты из точных координат графика, затем отдельно проверять каждый член при подстановке отрицательного аргумента.}
\itemhead{Исправление}
Сохраняем найденное $c=-2$ и правильно решаем уравнение для $a$:
\[
-2=4a-0{,}8\cdot2-2,\quad -2=4a-3{,}6,
\quad 4a=1{,}6,\quad a=0{,}4.
\]
\itemhead{Полное правильное решение}
На графике отмечены точки $(0;-2)$ и $(2;-2)$.
Первая точка даёт:
\[
f(0)=a\cdot0^2-0{,}8\cdot0+c=c=-2.
\]
Подставим вторую точку:
\[
-2=a\cdot2^2-0{,}8\cdot2-2
=4a-1{,}6-2.
\]
Отсюда $4a=1{,}6$, значит $a=0{,}4$.
Искомая функция:
\[
f(x)=0{,}4x^2-0{,}8x-2.
\]
Теперь вычисляем значение при $x=-5$:
\[
\begin{aligned}
f(-5)&=0{,}4\cdot(-5)^2-0{,}8\cdot(-5)-2\\
&=0{,}4\cdot25+4-2\\
&=10+4-2=12.
\end{aligned}
\]

\vishead{график найденной квадратичной функции}
\begin{center}
\begin{tikzpicture}[x=.82cm,y=.38cm]
 \begin{scope}
 \clip (-5.45,-3.45) rectangle (4.85,13.0);
 \draw[step=1,gray!22,very thin] (-5.45,-3.45) grid (4.85,13);
 \draw[-{Stealth},gray!75] (-5.45,0)--(4.65,0) node[right] {$x$};
 \draw[-{Stealth},gray!75] (0,-3.45)--(0,12.65) node[above] {$y$};
 \draw[blue!70!black,line width=1.4pt,samples=120,domain=-5.45:4.85,smooth,variable=\x]
  plot (\x,{.4*\x*\x-.8*\x-2});
 \end{scope}
 \fill[teal!75!black] (0,-2) circle (2.8pt) node[below right] {$(0;-2)$};
 \fill[teal!75!black] (2,-2) circle (2.8pt) node[below right] {$(2;-2)$};
 \fill[blue!70!black] (-5,12) circle (2.8pt) node[right] {$(-5;12)$};
 \draw[blue!70!black,dashed] (-5,0)--(-5,12)--(0,12);
 \node[below] at (-5,0) {$-5$}; \node[left] at (0,12) {$12$};
\end{tikzpicture}
\viscaption{Точки $(0;-2)$ и $(2;-2)$ задают коэффициенты; значение $f(-5)=12$ отмечено на той же параболе.}
\end{center}

\itemhead{Проверка}
Найденная функция принимает в узловых точках нужные значения: $f(0)=-2$ и $f(2)=1{,}6-1{,}6-2=-2$. Поэтому коэффициенты согласуются с графиком.
\conclusion{Искомое значение равно $12$.}
\itemhead{Задача на закрепление}
График функции $g(x)=ax^2-1{,}2x+c$ проходит через точки $(0;-3)$ и $(2;-1{,}4)$. Найдите $g(-4)$.

\task{14}{Отбор корней тригонометрического уравнения}
\itemhead{Текст задания}
а) Решите уравнение
\[
\cos^4\frac{x}{4}-\sin^4\frac{x}{4}
=\sin\left(x-\frac{\pi}{2}\right).
\]
б) Найдите все корни этого уравнения, принадлежащие отрезку $[\pi;5\pi]$.

\itemhead{Ваше решение}
В пункте (а) Вы верно раскрыли синус разности:
\[
\sin\left(x-\frac{\pi}{2}\right)=-\cos x.
\]
После разложения разности четвёртых степеней Вы правильно получили
\[
\cos\frac{x}{2}+\cos x=0,
\qquad 2\cos\frac{3x}{4}\cos\frac{x}{4}=0.
\]
Затем Вы записали обе серии корней:
\[
\begin{aligned}
 x&=\frac{2\pi}{3}+\frac{4\pi k}{3},&&k\in\mathbb Z,\\
 x&=2\pi+4\pi n,&&n\in\mathbb Z.
\end{aligned}
\]
Это корректное решение пункта (а). Отбора корней для пункта (б)
на предоставленной странице нет.

\itemhead{Ваш ответ}
Пункт (а): обе серии общих решений записаны верно.
Пункт (б): ответ не указан.

\err{Пункт (б) не завершён: не отобраны корни из найденных серий
на заданном отрезке $[\pi;5\pi]$. Ошибки в решении пункта (а) нет.}

\itemhead{Где именно возникает проблема}
Последний правильный шаг --- две общие серии корней,
приведённые выше. Первого неправильного математического перехода нет:
работа остановлена после выполнения пункта (а).
Необходимый следующий шаг --- ограничить целые параметры $k$ и $n$
условием принадлежности корня заданному отрезку.

\itemhead{Почему это неполное решение}
Запись $k,n\in\mathbb Z$ задаёт бесконечные семейства значений $x$.
Но пункт (б) требует только те значения, которые одновременно
принадлежат множеству решений уравнения и лежат в промежутке
от $\pi$ до $5\pi$, включая обе границы.

Здесь нельзя просто оставить общую формулу: нужно решить двойные
неравенства для параметров, перечислить допустимые целые значения
и убрать возможные повторы. В частности, корень $2\pi$
встречается сразу в обеих сериях, поэтому его следует записать один раз.

\keyidea{Начать с уже правильно найденных общих серий, подставить
каждую из них в двойное неравенство $\pi\le x\le5\pi$,
найти допустимые целые параметры и объединить полученные корни
без повторений.}

\itemhead{Исправление от последнего правильного шага}
Используем первую Вашу серию:
\[
 x=\frac{2\pi}{3}+\frac{4\pi k}{3},\qquad k\in\mathbb Z.
\]
Требуем, чтобы $\pi\le x\le5\pi$:
\[
\pi\le\frac{2\pi}{3}+\frac{4\pi k}{3}\le5\pi.
\]
Так как $\pi>0$, делим на $\pi$, затем умножаем на $3$:
\[
3\le2+4k\le15.
\]
Вычитаем $2$ и делим на $4$:
\[
1\le4k\le13,\qquad
\frac14\le k\le\frac{13}{4}.
\]
Следовательно, возможны только целые $k=1,2,3$.

\itemhead{Полное правильное решение}
Для полноты восстановим полученные Вами правильные преобразования
пункта (а):
\[
\begin{aligned}
\cos^4\frac{x}{4}-\sin^4\frac{x}{4}
&=\left(\cos^2\frac{x}{4}-\sin^2\frac{x}{4}\right)
\left(\cos^2\frac{x}{4}+\sin^2\frac{x}{4}\right)\\
&=\cos\frac{x}{2},\\
\sin\left(x-\frac\pi2\right)&=-\cos x.
\end{aligned}
\]
Поэтому
\[
\cos\frac{x}{2}+\cos x
=2\cos\frac{3x}{4}\cos\frac{x}{4}=0.
\]
Получаем уже найденные Вами серии
\[
\boxed{x=\frac{2\pi}{3}+\frac{4\pi k}{3}}
\quad\text{или}\quad
\boxed{x=2\pi+4\pi n},\qquad k,n\in\mathbb Z.
\]
Теперь выполним пункт (б).

\textbf{Первая серия.}
Из неравенства, разобранного выше, $k=1,2,3$.
Подставляем их по очереди:
\[
\begin{aligned}
 k=1&:\quad x=\frac{2\pi}{3}+\frac{4\pi}{3}=2\pi;\\
 k=2&:\quad x=\frac{2\pi}{3}+\frac{8\pi}{3}=\frac{10\pi}{3};\\
 k=3&:\quad x=\frac{2\pi}{3}+\frac{12\pi}{3}=\frac{14\pi}{3}.
\end{aligned}
\]

\textbf{Вторая серия.}
Рассмотрим Ваше выражение $x=2\pi+4\pi n$.
Условие принадлежности отрезку принимает вид
\[
\pi\le2\pi+4\pi n\le5\pi.
\]
Делим на $\pi>0$, затем вычитаем $2$:
\[
1\le2+4n\le5,
\qquad -1\le4n\le3,
\qquad -\frac14\le n\le\frac34.
\]
Здесь единственное целое значение параметра --- $n=0$.
Оно даёт $x=2\pi$.
Этот корень уже содержится в первой серии,
поэтому второй раз его не записываем.

Таким образом, все \emph{различные} корни на отрезке:
\[
\boxed{2\pi,\quad\frac{10\pi}{3},\quad\frac{14\pi}{3}}.
\]


\vishead{таблица отбора корней}
{\small\setlength{\tabcolsep}{4pt}
\begin{center}
\begin{tabular}{@{}p{31mm}p{53mm}p{31mm}p{40mm}@{}}
\toprule
\textbf{Серия} & \textbf{Выражение для $x$} & \textbf{Подходящие параметры} & \textbf{Корни на отрезке}\\
\midrule
Первая & $\dfrac{2\pi}{3}+\dfrac{4\pi k}{3}$ & $k=1,2,3$ & $2\pi,\ \dfrac{10\pi}{3},\ \dfrac{14\pi}{3}$\\[2pt]
Вторая & $2\pi+4\pi n$ & $n=0$ & $2\pi$ (повтор)\\
\bottomrule
\end{tabular}
\end{center}}
\smallskip
\begin{center}
\begin{tikzpicture}[x=1.75cm,y=1cm]
 \draw[-{Stealth},gray!80,thick] (.72,0)--(5.25,0) node[right] {$x$};
 \foreach \a in {1,2,3,4,5} {
   \draw[gray!65] (\a,-.08)--(\a,.08);
   \node[below] at (\a,-.1) {$\a\pi$};
 }
 \draw[blue!50!black,line width=2pt] (1,0)--(5,0);
 \foreach \a in {2,3.3333333,4.6666667} {\fill[teal!70!black] (\a,0) circle (2.6pt);}
 \node[above,teal!70!black] at (2,.16) {$2\pi$};
 \node[above,teal!70!black] at (3.3333,.16) {$\frac{10\pi}{3}$};
 \node[above,teal!70!black] at (4.6667,.16) {$\frac{14\pi}{3}$};
\end{tikzpicture}
\end{center}
\viscaption{Закрашены только различные корни, принадлежащие отрезку $[\pi;5\pi]$; значение $2\pi$ совпало в обеих сериях.}

\itemhead{Проверка}
Каждый из найденных корней принадлежит отрезку $[\pi;5\pi]$:
\[
\pi\le2\pi<\frac{10\pi}{3}<\frac{14\pi}{3}\le5\pi.
\]
Подставим их в преобразованное уравнение
$\cos(x/2)=-\cos x$:
\[
\begin{array}{c|cc}
 x&\cos(x/2)&-\cos x\\\hline
 2\pi&-1&-1\\
 10\pi/3&1/2&1/2\\
 14\pi/3&1/2&1/2
\end{array}
\]
Равенство выполняется в каждой строке.
Других целых $k,n$ в найденных диапазонах нет;
значит, все подходящие корни перечислены.

\conclusion{Пункт (а) выполнен верно. Для завершения пункта (б)
нужно было записать корни $2\pi$, $10\pi/3$ и $14\pi/3$.}

\itemhead{Задача на закрепление}
При решении тригонометрического уравнения получены
две серии корней:
\[
 x=\frac{\pi}{3}+\frac{2\pi k}{3},\quad k\in\mathbb Z;
 \qquad x=\pi+2\pi n,\quad n\in\mathbb Z.
\]
Найдите все различные корни на отрезке $[0;3\pi]$.

\task{16}{Показательное неравенство}
\itemhead{Текст задания}
Решите неравенство
\[
3\cdot25^{\,1-\frac3x}-152\cdot15^{-\frac3x}
+5\cdot3^{\,2-\frac6x}>0.
\]
\itemhead{Ваше решение}
Верно начато преобразование степеней по основаниям $5$ и $3$, затем введено обозначение $t=3/x$. Получено выражение с членами $75\cdot9^t$, $15^t$ и $45\cdot25^t$, но в последних строках коэффициент среднего члена записан как $159$ вместо исходного $152$. Ответ не получен.
\itemhead{Ваш ответ}
Не указан.
\err{При переписывании преобразованного выражения коэффициент $152$ изменён на $159$; кроме того, неравенство не решено до конца.}
\itemhead{Где именно возникает ошибка}
Последний надёжный правильный этап --- переход к $t=3/x$ и разложение $25=5^2$, $15=3\cdot5$, $3^2=9$.
Первый неверный переход виден в строке с общим знаменателем: вместо $-152\cdot15^t$ появляется $-159\cdot15^t$. Никакого действия, которое могло бы изменить $152$ на $159$, не выполнялось.
\itemhead{Почему это неверно}
Изменение числового коэффициента меняет соответствующее квадратное неравенство, его дискриминант, корни и итоговые интервалы. Также требуется учитывать, что $x\ne0$; обратная подстановка $t=3/x$ без контроля знака $x$ недопустима.
\keyidea{После умножения на положительную степень свести неравенство к квадратному относительно $z=(5/3)^{3/x}>0$ и аккуратно решить два неравенства для $3/x$.}
\itemhead{Исправление}
Сохраняем введённое $t=3/x$, но возвращаем исходный коэффициент $152$. Умножение обеих частей на $15^{2t}>0$ не меняет знак неравенства.
\itemhead{Полное правильное решение}
Область допустимых значений: $x\ne0$.
Пусть $t=3/x$. Тогда
\[
\begin{aligned}
3\cdot25^{1-t}&=75\cdot5^{-2t},\\
-152\cdot15^{-t}&=-152\cdot3^{-t}5^{-t},\\
5\cdot3^{2-2t}&=45\cdot3^{-2t}.
\end{aligned}
\]
Исходное неравенство принимает вид
\[
75\cdot5^{-2t}-152\cdot15^{-t}+45\cdot3^{-2t}>0.
\]
Умножим обе части на положительное число $15^{2t}$:
\[
75\cdot9^t-152\cdot15^t+45\cdot25^t>0.
\]
Теперь разделим на $9^t>0$:
\[
75-152\left(\frac53\right)^t
+45\left(\frac53\right)^{2t}>0.
\]
Обозначим $z=(5/3)^t>0$. Получаем квадратное неравенство
\[
45z^2-152z+75>0.
\]
Дискриминант:
\[
D=(-152)^2-4\cdot45\cdot75
=23104-13500=9604=98^2.
\]
Корни квадратного трёхчлена:
\[
z_1=\frac{152-98}{90}=\frac35,
\qquad z_2=\frac{152+98}{90}=\frac{25}{9}.
\]
Ветви параболы направлены вверх, поэтому
\[
z<\frac35\quad\text{или}\quad z>\frac{25}{9}.
\]
Так как $5/3>1$, показатель возрастает вместе со степенью. Заметим:
\[
\frac35=\left(\frac53\right)^{-1},\qquad
\frac{25}{9}=\left(\frac53\right)^2.
\]
Следовательно,
\[
\frac3x<-1\quad\text{или}\quad\frac3x>2.
\]
В первом случае правая часть отрицательна, поэтому $x<0$.
Умножая на отрицательное $x$, меняем знак:
\[
3>-x,\qquad x>-3.
\]
С учётом $x<0$ получаем $-3<x<0$.
Во втором случае $x>0$, и
\[
3>2x,\qquad 0<x<\frac32.
\]
Объединение интервалов:
\[
x\in(-3;0)\cup\left(0;\frac32\right).
\]

\vishead{знаки квадратного трёхчлена после замены}
\begin{center}
\begin{tikzpicture}[x=1.75cm,y=.9cm]
 \draw[-{Stealth},gray!75,thick] (-.08,0)--(4.05,0) node[right] {$z$};
 \draw[-{Stealth},gray!75] (0,-1.75)--(0,1.5) node[above] {$Q(z)$};
 \begin{scope}
 \clip (0,-1.73) rectangle (3.82,1.49);
 \draw[blue!70!black,line width=1.3pt,domain=0:3.9,samples=100,smooth,variable=\z]
   plot (\z,{.013*(45*\z*\z-152*\z+75)});
 \end{scope}
 \fill[teal!75!black] (.6,0) circle (2.4pt);
 \fill[teal!75!black] (2.77778,0) circle (2.4pt);
 \node[above left] at (.6,.02) {$\frac35$};
 \node[above right] at (2.77778,.02) {$\frac{25}{9}$};
 \node[font=\small,blue!65!black] at (1.7,-1.57) {$Q(z)<0$};
 \node[font=\small,teal!75!black] at (3.46,.85) {$Q(z)>0$};
\end{tikzpicture}
\viscaption{Здесь $Q(z)=45z^2-152z+75$. Парабола направлена вверх, поэтому $Q(z)>0$ вне корней; отдельно учитываем $z>0$.}
\end{center}
\vishead{перевод условий обратно к переменной $x$}
{\small\setlength{\tabcolsep}{5pt}
\begin{center}
\begin{tabular}{@{}p{40mm}p{47mm}p{65mm}@{}}
\toprule
\textbf{Условие для $z$} & \textbf{Неравенство для $3/x$} & \textbf{Промежуток $x$}\\
\midrule
$0<z<\frac35$ & $\frac3x<-1$ & $-3<x<0$\\[2pt]
$z>\frac{25}{9}$ & $\frac3x>2$ & $0<x<\frac32$\\
\bottomrule
\end{tabular}
\end{center}}

\itemhead{Проверка}
При $x=-1$ имеем $t=-3$, значит $z=(5/3)^{-3}<3/5$ и неравенство выполняется.
При $x=1$ получаем $t=3$ и $z=(5/3)^3>25/9$, неравенство выполняется.
При $x=3$ имеем $t=1$, то есть $z=5/3$ находится между корнями: неравенство не выполняется. Точка $x=0$ исключена из области определения.
\conclusion{Множество решений: $(-3;0)\cup(0;3/2)$.}
\itemhead{Задача на закрепление}
Решите неравенство
\[
3\cdot25^{\,1-\frac2x}-50\cdot15^{-\frac2x}
+5\cdot3^{\,2-\frac4x}>0.
\]

\task{17}{Минимизация расстояния по формуле}
\itemhead{Текст задания}
Расстояние вдоль противоположного берега реки вычисляется по формуле
\[
l=h\,\frac{u+v\cos\alpha}{v\sin\alpha}.
\]
Найдите наименьшее возможное $l$, если $h=75$ м, $u=2{,}6$ м/с, $v=1$ м/с, а угол направления лодки допускает переправу на противоположный берег.
\itemhead{Ваше решение}
Правильно получено
\[
l(\alpha)=\frac{195+75\cos\alpha}{\sin\alpha}.
\]
Но затем производная записана с переставленными членами числителя:
\[
l'_{\text{Ваш}}(\alpha)=
\frac{\cos\alpha(195+75\cos\alpha)
-\sin\alpha(-75\sin\alpha)}{\sin^2\alpha}.
\]
Из этого получено $\cos\alpha=-5/13$, $\sin\alpha=12/13$ и ответ $180$ м.
\itemhead{Ваш ответ}
$180$ м.
\err{В формуле производной частного перепутан порядок произведений. Числовой ответ оказался верным, но доказательство минимума не выполнено.}
\itemhead{Где именно возникает ошибка}
Последний верный шаг: $l(\alpha)=(195+75\cos\alpha)/\sin\alpha$.
Первый ошибочный шаг: вместо $u'v-uv'$ записано $uv'-u'v$.
Это меняет знак производной на противоположный. Условие равенства производной нулю при этом случайно остаётся тем же, поэтому найденное значение угла и итоговое число совпали с правильными.
\itemhead{Почему шаг неверен}
Если $l=u/v$, где $v\ne0$, то $l'=(u'v-uv')/v^2$.
Порядок вычитания принципиален: он определяет знак производной слева и справа от критической точки. Без корректного знака нельзя обосновать, что точка даёт минимум, а не максимум.
\keyidea{Сначала точно применить правило производной частного, затем проверить смену знака с минуса на плюс на физически допустимом промежутке $0<\alpha<\pi$.}

\vishead{модель переправы и сложения скоростей}
\begin{center}
\begin{tikzpicture}[scale=.81,>=Stealth]
  % Изображён оптимальный угол: cos(alpha)=-5/13, sin(alpha)=12/13;
  % результирующий путь имеет отношение l/h=180/75=2.4.
  \coordinate (A) at (2,0.75); \coordinate (B) at (2,3.75);
  \coordinate (C) at (9.2,3.75);
  \draw[gray!70,very thick] (.1,.75)--(10,.75);
  \draw[gray!70,very thick] (.1,3.75)--(10,3.75);
  \draw[dashed,gray!80] (A)--(B);
  \draw[teal!75!black,line width=1.3pt,-{Stealth}] (A)--(C)
    node[pos=.70,below right] {путь лодки};
  \draw[red!65!black,thick,-{Stealth}] (A)--(1.38,2.24)
    node[above left] {$\vec v$};
  \draw[blue!75!black,thick,-{Stealth}] (A)--(6.2,.75)
    node[pos=.62,below] {$\vec u$ (течение)};
  \draw[orange!65!black,thick,-{Stealth}] (2.0,0.75) ++(.71,0) arc[start angle=0,end angle=112.62,radius=.71];
  \node[orange!70!black] at (2.55,1.60) {$\alpha$};
  \draw[blue!55!black,line width=1.2pt] (B)--(C);
  \node[above] at (5.6,3.75) {$l=180\ \text{м}$};
  \node[right] at (2,2.55) {$h=75\ \text{м}$};
  \node[below left] at (A) {$A$};
  \node[above left] at (B) {$B$};
  \node[above right] at (C) {$C$};
  \draw (2,3.49)--(2.26,3.49)--(2.26,3.75);
\end{tikzpicture}
\viscaption{Схема для найденного оптимального угла: красная стрелка --- собственная скорость лодки $\vec v$, синяя --- течение $\vec u$. Угол $\alpha$ отсчитывается от направления течения, поэтому $\cos\alpha=-5/13<0$ и лодка направлена немного против течения.}
\end{center}

\itemhead{Исправление}
Пусть $u(\alpha)=195+75\cos\alpha$, $v(\alpha)=\sin\alpha$.
Тогда $u'(\alpha)=-75\sin\alpha$ и $v'(\alpha)=\cos\alpha$.
По правильному правилу:
\[
l'(\alpha)=\frac{(-75\sin\alpha)\sin\alpha
-(195+75\cos\alpha)\cos\alpha}{\sin^2\alpha}.
\]
\itemhead{Полное правильное решение}
Так как лодка пересекает реку, перпендикулярная берегу составляющая скорости должна быть положительна: $\sin\alpha>0$. Можно рассматривать $0<\alpha<\pi$.
Упростим числитель производной:
\[
\begin{aligned}
l'(\alpha)
&=\frac{-75\sin^2\alpha-195\cos\alpha-75\cos^2\alpha}{\sin^2\alpha}\\
&=-\frac{75(\sin^2\alpha+\cos^2\alpha)+195\cos\alpha}{\sin^2\alpha}\\
&=-\frac{75+195\cos\alpha}{\sin^2\alpha}.
\end{aligned}
\]
Приравниваем производную к нулю:
\[
75+195\cos\alpha=0,
\qquad \cos\alpha=-\frac{75}{195}=-\frac5{13}.
\]
На промежутке $0<\alpha<\pi$ синус положителен, поэтому
\[
\sin\alpha=\sqrt{1-\cos^2\alpha}
=\sqrt{1-\frac{25}{169}}=\frac{12}{13}.
\]
\textbf{Проверим минимум.}
При $\alpha$ меньше найденного угла функция $\cos\alpha$ больше $-5/13$. Тогда $75+195\cos\alpha>0$, поэтому $l'(\alpha)<0$.
При $\alpha$ больше найденного угла $\cos\alpha<-5/13$, значит $l'(\alpha)>0$.
Следовательно, функция сначала убывает, затем возрастает: найденная критическая точка даёт минимум.
Подставим значения синуса и косинуса:
\[
\begin{aligned}
l_{\min}
&=75\cdot\frac{\frac{13}{5}-\frac5{13}}{\frac{12}{13}}
=75\cdot\frac{\frac{169-25}{65}}{\frac{12}{13}}\\
&=75\cdot\frac{144}{65}\cdot\frac{13}{12}
=75\cdot\frac{12}{5}=180\ \text{м}.
\end{aligned}
\]

\vishead{знаки производной и минимум}
{\small\setlength{\tabcolsep}{5pt}
\begin{center}
\begin{tabular}{@{}p{40mm}p{42mm}p{31mm}p{37mm}@{}}
\toprule
\textbf{Угол} & \textbf{$75+195\cos\alpha$} & \textbf{$l'(\alpha)$} & \textbf{Поведение $l$}\\
\midrule
$0<\alpha<\alpha_0$ & $+$ & $-$ & убывает\\
$\alpha=\alpha_0$ & $0$ & $0$ & минимум\\
$\alpha_0<\alpha<\pi$ & $-$ & $+$ & возрастает\\
\bottomrule
\end{tabular}
\end{center}}
\viscaption{Здесь $\cos\alpha_0=-5/13$. Знак производной меняется с минуса на плюс --- это доказывает минимум.}

\itemhead{Проверка}
При приближении $\alpha$ к $0$ или к $\pi$ знаменатель $\sin\alpha$ стремится к нулю, а числитель остаётся положительным, поэтому $l$ неограниченно возрастает. Таким образом, найденный локальный минимум является наименьшим значением на всём допустимом промежутке.
\conclusion{Наименьшее расстояние равно $180$ м; числовой ответ совпал с Вашим, однако знак производной необходимо исправить.}
\itemhead{Задача на закрепление}
Расстояние до точки прибытия лодки выражается формулой
\[
L=60\,\frac{2+\cos\alpha}{\sin\alpha},
\qquad 0<\alpha<\pi.
\]
Найдите наименьшее возможное значение $L$.

\vfill
\begin{center}\small Лёвин Артём Александрович эксклюзивно для Николь Саркисьянц\end{center}
\end{document}